\subsection{Lineare Gleichungen}
\label{subsec:fundamentals_linear_equations}

Eine lineare Gleichung in einer Variablen hat die Form
\[
  ax+b=0, \qquad a\neq 0.
\]
\noindent\begin{minipage}{\linewidth}
  Durch Äquivalenzumformungen wird die Variable $x$ isoliert:
  \begin{ClozeArea}[27]{8}
    \[
      \begin{array}{r@{\quad}l@{\qquad}l}
        & ax+b=0 & \TransformOp{-b} \\[0.7em]
        \Longleftrightarrow & ax=-b
          & \TransformOp{:a \qquad (\text{nur für }a\neq0)} \\[0.7em]
        \Longleftrightarrow & x=-\dfrac{b}{a}. &
      \end{array}
    \]
  \end{ClozeArea}
\end{minipage}
Damit ist die Lösungsmenge
\[
  \mathcal{L}=\Cloze{5cm}{\left\{-\frac{b}{a}\right\}}.
\]
